% 202210-assessment-report.tex

% !TEX program = xelatex
% change it to xelatex if you use Chinese
\documentclass[]{beamer}

% preamble.tex

%%%%%%%%%%%%%%% Chinese %%%%%%%%%%%%%%%
%%%%% Uncomment the following lines to use Chinese
\usepackage{xeCJK}
\usepackage{fontspec}
%%%%%%%%%%%%%%% Chinese %%%%%%%%%%%%%%%

%%%%%%%%%%%%%%% Theme %%%%%%%%%%%%%%%
\usetheme{CambridgeUS} % try Madrid
\usecolortheme{beaver} % try beaver, dolphin, seahorse
% \usefonttheme[onlymath]{serif} % try "professionalfonts"
\usefonttheme{serif}  % standard font (same with that in ``standalone'')
% \setCJKmainfont{Microsoft YaHei} % try SimSun
% \setmainfont{Gill Sans}
% \setsansfont{Gill Sans}

\usepackage[export]{adjustbox}

\setbeamersize{text margin left = 2em, text margin right = 1em}
\setbeamercolor{footnote mark}{fg = teal}
\setbeamertemplate{itemize items}[default]
\setbeamertemplate{enumerate items}[default]

\renewcommand*{\thefootnote}{\alph{footnote}}
%%%%%%%%%%%%%%% Theme %%%%%%%%%%%%%%%

\usepackage{graphicx, subcaption}

\usepackage{amssymb, pifont}
\newcommand{\cmark}{\ding{51}}
\newcommand{\xmark}{\ding{55}}

%%%%%%%%%%%%%%% Algorithms %%%%%%%%%%%%%%%
\usepackage{algorithm}
\usepackage[noend]{algpseudocode}
% for two-column algorithms
% see https://tex.stackexchange.com/a/159431.
\usepackage{varwidth}

\newcommand{\code}[2]{#1:#2}
\newcommand{\hStatex}[0]{\vspace{4pt}}
\renewcommand\algorithmicthen{}
\renewcommand\algorithmicdo{}

\algblockdefx[upon]{When}{EndWhen}%
  [3][Unknown]{\textbf{when for every} #3 \textbf{received a quorum of} #1(#2)}%
  {\textbf{end}}

\algblockdefx[upon]{WhenReceive}{EndWhenReceive}%
  [1][Unknown]{\textbf{when received} #1}%
  {\textbf{end}}

\algblockdefx[upon]{Upon}{EndUpon}%
  [1][Unknown]{\textbf{upon} #1}%
  {\textbf{end upon}}

% Remove the end blank line when using noend in conjunction with custom blocks
% See https://tex.stackexchange.com/a/172065
\makeatletter
\ifthenelse{\equal{\ALG@noend}{t}}%
  {\algtext*{EndWhen}}
  {}%
\makeatother

\makeatletter
\ifthenelse{\equal{\ALG@noend}{t}}%
  {\algtext*{EndWhenReceive}}
  {}%
\makeatother

\makeatletter
\ifthenelse{\equal{\ALG@noend}{t}}%
  {\algtext*{EndUpon}}
  {}%
\makeatother
%%%%%%%%%%%%%%% Algorithms %%%%%%%%%%%%%%%

%%%%%%%%%%%%%%% Theorems %%%%%%%%%%%%%%%
% theorems (global numbering)
\theoremstyle{definition}
\newtheorem{property}[theorem]{Property}
\newtheorem{assumption}{\textsc{Assumption}}
%%%%%%%%%%%%%%% Theorems %%%%%%%%%%%%%%%

\usepackage{caption}
\DeclareCaptionLabelSeparator{none}{}
\captionsetup{labelsep = none}

\makeatletter
\let\@@magyar@captionfix\relax
\makeatother

%%%%%%%%%%%%%%% Tables %%%%%%%%%%%%%%%
\usepackage{multirow}
\newcommand{\incell}[2]{\begin{tabular}[c]{@{}c@{}}#1\\[3pt] #2\end{tabular}}
%%%%%%%%%%%%%%% Tables %%%%%%%%%%%%%%%

%%%%%%%%%%%%%%% Figures %%%%%%%%%%%%%%%
% for fig without caption: #1: width/size; #2: fig file
\newcommand{\fig}[2]{
  \begin{figure}[htp]
    \centering
    \includegraphics[#1]{#2}
  \end{figure}
}

% for fig with caption: #1: width/size; #2: fig file; #3: fig caption
\newcommand{\figcap}[3]{
  \begin{figure}[htp]
    \centering
    \includegraphics[#1]{#2}
    \caption{#3}
  \end{figure}
}

%%%%% For Chinese %%%%%
% \renewcommand\figurename{图:\;}
% \renewcommand\tablename{表:\;}
%%%%% For Chinese %%%%%
%%%%%%%%%%%%%%% Figures %%%%%%%%%%%%%%%

%%%%%%%%%%%%%%% TiKZ %%%%%%%%%%%%%%%
\usepackage[linewidth = 1pt, framemethod = TikZ]{mdframed}
\mdfsetup{frametitlealignment=\center}
%%%%%%%%%%%%%%% TiKZ %%%%%%%%%%%%%%%

%%%%%%%%%%%%%%% Citation %%%%%%%%%%%%%%%
% try "backend = bibtex", "backend = biber"
\usepackage[natbib = true, backend = biber, style = authoryear, maxbibnames = 99]{biblatex}
\setbeamertemplate{bibliography item}[article]
\renewcommand*{\bibfont}{\footnotesize}
\addbibresource{2022-assessment-report.bib}

\newcommand{\ncite}[1]{\violet{\footnotesize [\cite{#1}]}}
%%%%%%%%%%%%%%% Citation %%%%%%%%%%%%%%%

%%%%%%%%%%%%%%% Colors %%%%%%%%%%%%%%%
\usepackage{xcolor}

\definecolor{DarkRed}{rgb}{0.55, 0.0, 0.0}

\newcommand{\red}[1]{\textcolor{red}{#1}}
\newcommand{\green}[1]{\textcolor{green}{#1}}
\newcommand{\blue}[1]{\textcolor{blue}{#1}}
\newcommand{\purple}[1]{\textcolor{purple}{#1}}
\newcommand{\cyan}[1]{\textcolor{cyan}{#1}}
\newcommand{\violet}[1]{\textcolor{violet}{#1}}
\newcommand{\lgray}[1]{\textcolor{lightgray}{#1}}
\newcommand{\teal}[1]{\textcolor{teal}{#1}}
\newcommand{\brown}[1]{\textcolor{brown}{#1}}
\newcommand{\orange}[1]{\textcolor{orange}{#1}}

\newcommand{\hl}[2]{\fcolorbox{#1}{#1!60}{#2}}
%%%%%%%%%%%%%%% Colors %%%%%%%%%%%%%%%

%%%%%%%%%%%%%%% thankyou %%%%%%%%%%%%%%%
\newcommand{\thankyou}{
  \begin{frame}[noframenumbering]
    \begin{center}
      \fig{width = 0.618\textwidth}{figs/thankyou}
      \vspace{0.30cm}
      Hengfeng Wei (hfwei@nju.edu.cn)
    \end{center}
  \end{frame}
}
%%%%%%%%%%%%%%% thankyou %%%%%%%%%%%%%%%
%%%%%%%%%%%%%%%%%%%%%%%%%%%%%%%%%%%%%%%%%%%%%%%%%%%%%%%%%%%%%%%%%%%%%%%%%%%%%%%%
\title[]{分布数据一致性: 形式化规约与验证}
\subtitle{(软件学院聘期考核报告; 2020年01月 $\sim$ 2022年10月)}

\author[魏恒峰]{魏恒峰}
\titlegraphic{\includegraphics[height = 1.5cm]{figs/nju-logo-purple}~
              \includegraphics[height = 1.5cm]{figs/software-logo}}
\institute{联系教授: 马晓星 (软件所)}
\date{2022年11月04日}
%%%%%%%%%%%%%%%%%%%%
\begin{document}
\maketitle
%%%%%%%%%%%%%%%%%%%%

% teaching.tex

%%%%%%%%%%%%%%%%%%%%
\begin{frame}{}
  \fig{width = 0.80\textwidth}{figs/teach}
\end{frame}
%%%%%%%%%%%%%%%%%%%%

%%%%%%%%%%%%%%%%%%%%
\begin{frame}{}
  \begin{table}[t]
    \centering
    \renewcommand\arraystretch{1.2}
    \begin{tabular}[]{c|c|c|c}
      \hline
      学期       & 课程                & 学分    & 课时 \\ \hline \hline
      2023年春季 & 编译原理 (1 班)           & 3  & 54    \\ \hline
      2023年暑期 & 大语言模型原理与应用        & 1  & 2    \\ \hline
      2023年秋季 & C 语言程序设计基础 (1 班)   & 2  & 36    \\ \hline
      2023年秋季 & C 语言程序设计基础 (2 班)   & 2  & 36    \\ \hline
    \end{tabular}
  \end{table}

  \begin{columns}
    \column{0.33\textwidth}
      \fig{width = 0.40\textwidth}{figs/cpl-logo}
    \column{0.35\textwidth}
      \fig{width = 0.50\textwidth}{figs/compiler-logo}
  \end{columns}
\end{frame}
%%%%%%%%%%%%%%%%%%%%

%%%%%%%%%%%%%%%%%%%%
\begin{frame}{}
  \[
    (\underbrace{200}_{\text{\red{软件学院}}}
      + \underbrace{88}_{\text{\red{跨专业选修}}})
      + (\underbrace{98 + 95 + 87 + 90 + 86 + 90}_{\text{技术科学试验班}})
      + \underbrace{32}_{\text{苏州校区重修班}} = 866 \text{ 名学生}
  \]

  \fig{width = 0.30\textwidth}{figs/cpl-logo}

  \[
    (\underbrace{8}_{\text{\scriptsize 软件学院}}
      + (\underbrace{3 \times 6}_{\text{技术科学试验班}})
      + \underbrace{1}_{\text{苏州校区重修班}} = 27 \text{ 名助教}
  \]
\end{frame}
%%%%%%%%%%%%%%%%%%%%

%%%%%%%%%%%%%%%%%%%%
\begin{frame}{}
  \begin{center}
    10 月 29 日, 已顺利完成第一次机考

    \fig{width = 0.95\textwidth}{figs/cpl-scores}

    \vspace{0.60cm}
    定于 12 月 09 日, 第二次机考
  \end{center}
\end{frame}
%%%%%%%%%%%%%%%%%%%%

%%%%%%%%%%%%%%%%%%%%
\begin{frame}{}
  \begin{center}
    每周安排 9 次答疑
    \fig{width = 0.70\textwidth}{figs/cpl-qa}
  \end{center}
\end{frame}
%%%%%%%%%%%%%%%%%%%%

%%%%%%%%%%%%%%%%%%%%
\begin{frame}{}
  \begin{center}
    答疑收集表
    \fig{width = 1.00\textwidth}{figs/qa-table}
  \end{center}
\end{frame}
%%%%%%%%%%%%%%%%%%%%

%%%%%%%%%%%%%%%%%%%%
\begin{frame}{}
  \fig{width = 1.00\textwidth}{figs/cpl-bilibili-lectures}
\end{frame}
%%%%%%%%%%%%%%%%%%%%

%%%%%%%%%%%%%%%%%%%%
\begin{frame}{}
  \fig{width = 1.00\textwidth}{figs/cpl-bilibili-outofclass}
\end{frame}
%%%%%%%%%%%%%%%%%%%%

%%%%%%%%%%%%%%%%%%%%
\begin{frame}{}
  \begin{center}
    2023 春季, 《编译原理》由选修课改为专业必修课。

    \vspace{0.30cm}
    \fig{width = 0.40\textwidth}{figs/compiler-logo}
  \end{center}
\end{frame}
%%%%%%%%%%%%%%%%%%%%

%%%%%%%%%%%%%%%%%%%%
\begin{frame}{}
  \begin{center}
    本学期: \cyan{\bf 作业 (0 分) + \purple{实验 (75 分)} + 期末测试 (25 分)}

    \vspace{0.80cm}

    \cyan{\bf 作业 (15 分) + \purple{实验 (45 分)} + 期末测试 (40 分)}

    \vspace{1.50cm}
    实验分数高, 导致今年的高分段人数较多。

    \vspace{1.00cm}
    下学期考虑调整。
  \end{center}
\end{frame}
%%%%%%%%%%%%%%%%%%%%

%%%%%%%%%%%%%%%%%%%%
\begin{frame}{}
  \begin{center}
    \fig{width = 0.50\textwidth}{figs/llvm-riscv}

    \vspace{0.30cm}
    本学期引入了 RISC-V, 为``目标代码生成阶段''做准备
  \end{center}
\end{frame}
%%%%%%%%%%%%%%%%%%%%

%%%%%%%%%%%%%%%%%%%%
\begin{frame}{}
  \begin{center}
    计划编写《编译原理》课程讲义

    \vspace{0.80cm}
    逐步对外开放《C 语言程序设计基础》与《编译原理》课程资源

    \vspace{0.30cm}
    提升课程影响力
  \end{center}
\end{frame}
%%%%%%%%%%%%%%%%%%%%
% research.tex

%%%%%%%%%%%%%%%%%%%%
\begin{frame}{}
	\fig{width = 0.80\textwidth}{figs/research}
\end{frame}
%%%%%%%%%%%%%%%%%%%%

%%%%%%%%%%%%%%%%%%%%
\begin{frame}{}
	\fig{width = 1.00\textwidth}{figs/VLDB}
	\fig{width = 1.00\textwidth}{figs/SIGMOD}
	\fig{width = 0.90\textwidth}{figs/TPDS}
\end{frame}
%%%%%%%%%%%%%%%%%%%%

%%%%%%%%%%%%%%%%%%%%
\begin{frame}{}
	\begin{center}
		数据库系统与客户程序的正确性是至关重要的

		\fig{width = 0.65\textwidth}{figs/problems}

		{\small 三个事务一致性验证问题之间的递进关系}
	\end{center}
\end{frame}
%%%%%%%%%%%%%%%%%%%%

%%%%%%%%%%%%%%%%%%%%
\begin{frame}{}
  \begin{center}
    \blue{\bf 我们的贡献:} 高效的快照隔离检测算法与工具~\ncite{PolySI:VLDB2023}

    \fig{width = 0.85\textwidth}{figs/polysi-vldb2023}
  \end{center}
\end{frame}
%%%%%%%%%%%%%%%%%%%%

%%%%%%%%%%%%%%%%%%%%
\begin{frame}{}
  \begin{center}
    \fig{width = 0.70\textwidth}{figs/checker-polysi}
  \end{center}
\end{frame}
%%%%%%%%%%%%%%%%%%%%

%%%%%%%%%%%%%%%%%%%%
\begin{frame}{}
	\centerline{性能显著优于其它 SI 检测算法 (在腾讯 TDSQL 开发过程初步测试)}

	\fig{width = 0.60\textwidth}{figs/polysi-runtime}
  %\#sessions=20, \#txns/session=100, \#ops/txn=15,   keys=10k,  \%read=50\%, distribution=zipfian.
\end{frame}
%%%%%%%%%%%%%%%%%%%%

%%%%%%%%%%%%%%%%%%%%%%%%%%%%%%
\begin{frame}{}
	\begin{center}
		\polysi{} 以黑盒的方式使用 MonoSAT 求解器

		\vspace{0.30cm}
		\fig{width = 0.70\textwidth}{figs/solver-blackbox}
		\vspace{0.30cm}

		未能充分利用 SAT/SMT 搜索框架 (DPLL/CDCL)
	\end{center}
\end{frame}
%%%%%%%%%%%%%%%%%%%%%%%%%%%%%%

%%%%%%%%%%%%%%%%%%%%%%%%%%%%%%
\begin{frame}{}
	\begin{center}
		\blue{\bf 当前的探索:} 如何``深度整合 SMT'', 进一步提升算法效率?

		\vspace{0.50cm}
		\fig{width = 0.70\textwidth}{figs/integrated}
	\end{center}
\end{frame}
%%%%%%%%%%%%%%%%%%%%%%%%%%%%%%

%%%%%%%%%%%%%%%%%%%%%%%%%%%%%%
\begin{frame}{}
  \begin{center}
		\red{\bf 进行中:} 设计专用的事务一致性理论求解器, 初见成效

	  \fig{width = 0.95\textwidth}{figs/consistency-theory-solver-revised}
  \end{center}
\end{frame}
%%%%%%%%%%%%%%%%%%%%%%%%%%%%%%

%%%%%%%%%%%%%%%%%%%%
\begin{frame}{}
	\begin{center}
		``基于深度整合 SMT 方案的事务一致性高效验证技术''
		\fig{width = 1.00\textwidth}{figs/roadmap}
	\end{center}
\end{frame}
%%%%%%%%%%%%%%%%%%%%

%%%%%%%%%%%%%%%%%%%%
\begin{frame}{}
	\begin{columns}
		\column{0.60\textwidth}
		  \fig{width = 1.00\textwidth}{figs/general-program-rejected}
		\column{0.40\textwidth}
		  \begin{center}
				去年已结题

				\fig{width = 1.00\textwidth}{figs/tencent-UR-logo}

				\vspace{0.30cm}
				继续与腾讯团队保持合作 (暂无项目支持)
			\end{center}
	\end{columns}

	\vspace{0.30cm}
	\begin{center}
	  \red{\bf 亟需加强项目申报}
	\end{center}
\end{frame}
%%%%%%%%%%%%%%%%%%%%

%%%%%%%%%%%%%%%%%%%%
\begin{frame}{总结}
  \begin{center}
    魏恒峰 (hfwei@nju.edu.cn)
  \end{center}

  \begin{table}[]
  \centering
    \resizebox{0.85\textwidth}{!}{%
      \begin{tabular}{|c||c|c|c|}
      \hline
	{\bf 聘期合同要求} 		& {\bf 工作情况}  \\ \hline\hline
	\red{\bf 教学:} 承担一门课程    		& \incell{问题求解课程}{\incell{五个学期; 共164学时}{(2019级本科生``我心目中的好课程'')}} \\ \hline
	\red{\bf 科研:} 4-6篇高水平论文		& \incell{发表3篇 (含1篇短文)}{\incell{在审4篇}{(2017年CCF优秀博士学位论文奖)}}  \\ \hline
	人才培养		& \incell{负责或协助指导学生9人次}{(学术积累: 组织TLA$^{+}$与Coq讨论班)}  \\ \hline
	\incell{主持/参与}{多个基金项目}	& \incell{主持1项; 参与1项}{个人可支配总经费75万元} \\ \hline
      \end{tabular}%
    }
  \end{table}
\end{frame}
%%%%%%%%%%%%%%%%%%%%

%%%%%%%%%%%%%%%%%%%%
\thankyou{}
%%%%%%%%%%%%%%%%%%%%
\begin{frame}[allowframebreaks]
  \printbibliography
\end{frame}
%%%%%%%%%%%%%%%%%%%%
\end{document}
%%%%%%%%%%%%%%%%%%%%